\daytitle{1}{Pure fixed points of a channel}

\begin{daymeta}
  \item[Topic] Quantum channels; Kraus operators; positive semidefinite arguments.
  \item[Estimated time] $\sim 12$ minutes (solving and write-up).
  \item[Origin] Original. The underlying facts are folklore in the theory of fixed
    points of CPTP maps; the formulation and the two-part structure are mine.
  \item[Source] Original problem.
  \item[Prerequisites] Kraus representation of a channel.
\end{daymeta}

\begin{problem}
Let $\Phi$ be a quantum channel (CPTP map) on $\Lop{\CC^{d}}$, with Kraus decomposition
\[
  \Phi(X)=\sum_{i} K_i X K_i^{\dagger},
  \qquad
  \sum_i K_i^{\dagger} K_i = \id .
\]
Suppose $\Phi$ has a \emph{pure} fixed point: there is a unit vector
$\ket{\psi}\in\CC^{d}$ with $\Phi(\proj{\psi})=\proj{\psi}$.
\begin{enumerate}[label=\textbf{(\alph*)},leftmargin=*]
  \item Show that every Kraus operator satisfies $K_i\ket{\psi}=c_i\ket{\psi}$ for some
        $c_i\in\CC$, and determine the constraint on the $c_i$. Show also that this
        conclusion does not depend on which Kraus decomposition of $\Phi$ was chosen.
  \item Prove that $\Phi$ can never decrease the overlap with $\ket{\psi}$: for
        \emph{every} state $\sigma\in\Dop{\CC^{d}}$,
        \[
          \avg{\psi}{\Phi(\sigma)}{\psi} \;\ge\; \avg{\psi}{\sigma}{\psi} .
        \]
\end{enumerate}
\end{problem}

% ============================================================== solutions ===
\subsection*{Solution}

\begin{solution}[Part (a)]
Define the orthonormal basis $\ket{k}$, $k=0,...,n-1$, with $\ket{0}=\ket{\psi}$, we than have that for any $k>0$ 
\begin{equation}
    \bra{k}\Phi(\proj{0})\ket{k} = \sum_i |\bra{k}K_i\ket{0}|^2 =0
\end{equation}
Thus $\bra{k}K_i\ket{0}=0, \forall k>0$, implying that $K_i\ket{0}$ can only have component in the $\ket{0}$ direction, i.e. there is $c_i \in \CC$ such that $K_i\ket{0}=c_i\ket{0}$, which must satisfy
\begin{equation}
    1 = \avg{0}{\sum_i K_i^{\dagger} K_i}{0} = \sum_i |c_i|^2
\end{equation}

\end{solution}

\begin{solution}[Part (b)]
\begin{equation}
    \avg{0}{\Phi(\sigma)}{0} = \Tr(\proj{0}\Phi(\sigma))=\Tr(\sum_i\proj{0}K_i\sigma K_i^\dagger)=\Tr(\sum_iK_i^\dagger\proj{0}K_i\sigma )
\end{equation}
Define $A=\sum_iK_i^\dagger\proj{0}K_i$, thus $A\geq 0$ directly from its expression and $A \leq \id$ since $ \sum_i K_i^{\dagger} K_i = \id$. 

Now, we have that $\avg{0}{A}{0}=\sum_i |c_i|^2=1$ combining this with $A\leq 1$, we get that this is only possible if $(\id - A) \ket{0}=0$, so $A\geq \proj{0}$, i.e. $A-\proj{0} \geq 0$. 

But $\sigma \geq 0$ as well, so $\Tr((A-\proj{0})\sigma) \geq 0$, therefore
\begin{equation}
\avg{\psi}{\Phi(\sigma)}{\psi}= \Tr(A\sigma) \;\ge\; \Tr(\proj{0}\sigma)=\avg{\psi}{\sigma}{\psi}
\end{equation}
\end{solution}



% =============================================================== bookkeeping =
\begin{attemptlog}
  \item[Date] 04/09/2026
  \item[Time taken] 45 min
  \item[Difficulty] 3
  \item[User Opinion] successful insight
\end{attemptlog}